\input{/Users/tommasopedroni/Desktop/TESI/Tesi/latex_prima_versione/preambolo_bozza.tex}
\newgeometry{margin=2cm}
\setlength{\headwidth}{\textwidth}

\begin{document}
\setcounter{chapter}{3}
\chapter{Von Neumann Algebras and The Modular Operator}
\section{Von Neumann Algebras, Concretely and Abstractly}\label{sec:vn-algebras}

The previous chapter derived Schr\"odinger dynamics, the Born rule, and
Noether's theorem from the intrinsic K\"ahler geometry of the projective
Hilbert space, once a state and a distinguished selfadjoint Hamiltonian
were given. The role of the Hamiltonian itself, however, was never
questioned. In every construction so far, a time evolution enters the
theory as external structure, whether as the one-parameter group of
Definition 1.6.23 in the classical setting or as the Hamiltonian flow of
Proposition 3.3.2 in the quantum one. This chapter asks whether that
structure can instead be recovered from the algebra together with a
single state, without postulating a Hamiltonian at all. The
Tomita--Takesaki theorem shows that every faithful normal state on a von
Neumann algebra determines a canonical one-parameter group of
automorphisms, its modular automorphism group. When the state describes
thermal equilibrium, this group coincides with ordinary time evolution
up to a rescaling by the inverse temperature, an identification that
suggests reading the modular flow itself as the origin of physical time
whenever no external clock is available. This is the content of the
thermal time hypothesis of Connes and Rovelli, and it is the
algebraic-geometric standpoint developed throughout this thesis that
makes the proposal precise.

Every construction so far in this thesis singled out the $C^*$-algebra
of Definition 1.6.11 as the primitive algebraic object, and derived a
Hilbert space representation of an observable algebra only afterward,
through the Gelfand-Naimark-Segal construction of Theorem 2.7. The
states considered there were arbitrary positive normalized functionals,
with no continuity required beyond what automatic continuity of
Proposition 2.1.23 already supplies. Tomita--Takesaki theory, together
with the notion of a thermal state occupying Section 4.5, requires a
strictly stronger setting. The relevant algebra of observables must be
weakly closed, so that it coincides with its own bicommutant, and the
state that generates the modular flow must be normal, so that
expectation values remain continuous under the monotone limits used to
define equilibrium. This section introduces von Neumann algebras and
normal states as this strengthened setting, and isolates the property
under which the Gelfand-Naimark-Segal vector of Theorem 2.7 is
simultaneously cyclic for the algebra and separating for it, a
distinction the construction of Section 4.3 will use directly. A
concrete algebra of operators on a fixed Hilbert space, however, is not
the natural datum on the abstract side of the correspondence developed
in Chapters 1 and 2, where the observable algebra $A$ of Definition
2.1.22 carries no reference to any particular Hilbert space. The second
half of this section therefore isolates the property of $A$ alone,
having a predual, that identifies exactly when $A$ can be realized this
way, and closes a gap left open by the Gelfand-Naimark-Segal
construction itself: nothing in Theorem 2.7 guaranteed that the
represented algebra $\pi_\omega(A)$ is weakly closed, hence a genuine
von Neumann algebra in the sense about to be introduced.

\begin{definition}[Commutant]\label{def:commutant}
	Let $H$ be a Hilbert space, as per Definition 2.2.6, and let
	$S \subseteq B(H)$ be a subset. The commutant of $S$ is
	\[
	S' := \{T \in B(H) \mid TA = AT \ \forall A \in S\}.
	\]
\end{definition}

\begin{remark}
	The commutant $S'$ is a subalgebra of $B(H)$ containing the identity,
	and every element of $S$ commutes with every element of $S'$ by
	definition, so $S$ is contained in its double commutant $S'' := (S')'$.
\end{remark}

The characterization of an algebra through its double commutant, rather
than through an explicit topological closure, is only useful once the
relevant topology on $B(H)$ has itself been fixed. The operator norm
topology used throughout this thesis is too fine for this purpose,
since it does not detect the monotone limits under which a state is
required to behave continuously in Definition~\ref{def:normal-state}
below, so a coarser pair of topologies is introduced first.

\begin{definition}[Weak and strong operator topologies]\label{def:wot-sot}
	Let $H$ be a Hilbert space, as per Definition 2.2.6. The weak operator
	topology on $B(H)$ is the coarsest topology for which every map
	$T \mapsto (\phi, T\psi)$, with $\phi, \psi \in H$, is continuous. The
	strong operator topology is the coarsest topology for which every map
	$T \mapsto \lVert T\psi \rVert$, with $\psi \in H$, is continuous.
\end{definition}

With both topologies available, it becomes possible to ask whether the
algebraic condition $M = M''$ used above coincides with topological
closure. Von Neumann's original theorem answers this affirmatively for
unital $*$-subalgebras, and this equivalence is what justifies treating
the double commutant, an entirely algebraic condition, as the working
definition of the topologically closed algebras this chapter needs.

\begin{theorem}[von Neumann Bicommutant Theorem]\label{thm:bicommutant}
	Let $M \subseteq B(H)$ be a unital $*$-subalgebra, as per
	Definition 1.6.2. The following are equivalent.
	\begin{enumerate}
		\item $M = M''$, where $M''$ denotes the double commutant of
		Definition~\ref{def:commutant}.
		\item $M$ is closed in the weak operator topology of
		Definition~\ref{def:wot-sot}.
		\item $M$ is closed in the strong operator topology of
		Definition~\ref{def:wot-sot}.
	\end{enumerate}
\end{theorem}
\noindent This is proved in \cite[TODO: collocazione precisa]{TODO-vonNeumannBicommutant}.

\begin{definition}[von Neumann algebra]\label{def:vN-algebra}
	A von Neumann algebra on $H$ is a unital $*$-subalgebra
	$M \subseteq B(H)$, as per Definition 1.6.2, satisfying $M = M''$, as
	per Definition~\ref{def:commutant}.
\end{definition}

\begin{remark}
	By Theorem~\ref{thm:bicommutant}, a von Neumann algebra can equally be
	characterized as a unital $*$-subalgebra of $B(H)$ closed in either
	operator topology of Definition~\ref{def:wot-sot}. Every von Neumann
	algebra is in particular a $C^*$-algebra, as per Definition 1.6.11,
	under the operator norm, since the norm topology is finer than both
	operator topologies and closure therefore persists from either of them
	to the norm topology.
\end{remark}

The states relevant to modular theory must interact well with the
operator topologies of Definition~\ref{def:wot-sot}, not merely with
the norm topology already available on any $C^*$-algebra by Definition
2.1.3. The most direct way to express this compatibility is to exhibit
each such state as the restriction of a single linear functional
defined on the whole of $B(H)$, built from a fixed trace-class operator
rather than from the algebra $M$ itself. This requires first extending
the trace of Definition 2.3.15, so far used only for finite-rank
density matrices in Proposition 2.3.14, to a linear functional defined
on a genuine two-sided ideal of $B(H)$.

\begin{definition}[Trace of a positive operator]\label{def:trace}
	Let $H$ be a Hilbert space, as per Definition 2.2.6, and let
	$\{e_i\}_{i \in I}$ be an orthonormal basis of $H$. For $A \in B(H)$
	positive, as per Definition 2.2.22, the trace of $A$ is
	\[
	\mathrm{Tr}(A) := \sum_{i \in I} (e_i, A e_i) \in [0, \infty].
	\]
\end{definition}

\begin{remark}
	The value of $\mathrm{Tr}(A)$ in Definition~\ref{def:trace} does not
	depend on the choice of orthonormal basis $\{e_i\}_{i\in I}$, so that
	$\mathrm{Tr}$ is a well defined map from the positive elements of
	$B(H)$ to $[0,\infty]$. This invariance is proved in \cite[TODO:
	collocazione precisa]{TODO-traceclass}.
\end{remark}

Definition~\ref{def:trace} assigns a value in $[0,\infty]$ to every
positive operator, but a genuine linear functional requires this value
to be finite everywhere, and a genuine ideal of $B(H)$ requires closure
not only under sums of positive elements but under multiplication by an
arbitrary bounded operator on either side. Both properties hold once
the trace is imposed on the modulus of a general, not necessarily
positive, operator.

\begin{definition}[Trace-class operators]\label{def:trace-class}
	Let $H$ be a Hilbert space, as per Definition 2.2.6, and let
	$A \in B(H)$. Write $|A| := (A^*A)^{1/2}$, the positive square root of
	$A^*A$ as per Remark 2.2.23, positive since
	$(\psi, A^*A\psi) = \lVert A\psi \rVert^2 \geq 0$ for every $\psi \in
	H$. The operator $A$ is trace-class if $\mathrm{Tr}(|A|) < \infty$, as
	per Definition~\ref{def:trace}. The set of trace-class operators on $H$
	is denoted $B_1(H)$.
\end{definition}

\begin{remark}
	The notation $B_1(H)$ parallels the notation $B_0(H)$ of Definition
	2.2.15 for the ideal of compact operators. In fact, $B_1(H)$ is itself
	a two-sided ideal of $B(H)$, contained in $B_0(H)$, and the trace
	extends from the positive elements of $B_1(H)$ to a linear functional
	on the whole of $B_1(H)$, satisfying
	\[
	\mathrm{Tr}(AB) = \mathrm{Tr}(BA), \qquad A \in B_1(H), \ B \in B(H).
	\]
	Setting $\lVert A \rVert_1 := \mathrm{Tr}(|A|)$ makes
	$(B_1(H), \lVert \cdot \rVert_1)$ a Banach space, as per Definition
	1.6.9, and the map
	\[
	B(H) \to (B_1(H))', \qquad T \mapsto \big( S \mapsto \mathrm{Tr}(ST)
	\big),
	\]
	with $(B_1(H))'$ the topological dual of Definition 2.2.13, is an
	isometric isomorphism.\label{eq:trace-duality-remark} Every fact
	recalled in this remark is proved in \cite[TODO: collocazione
	precisa]{TODO-traceclass}.
\end{remark}

The isometric identification of $B(H)$ with the topological dual of
$B_1(H)$ singles out a topology on $B(H)$ coarser than the operator
norm, and, as the following remark records, also finer than the two
topologies of Definition~\ref{def:wot-sot} in general, yet coincident
with them on the bounded sets relevant to representing states.

\begin{definition}[Ultraweak topology]\label{def:ultraweak}
	Let $H$ be a Hilbert space, as per Definition 2.2.6. The ultraweak
	topology on $B(H)$ is the coarsest topology for which every evaluation
	map $\widehat{S} : B(H) \to \mathbb{C}$, $\widehat{S}(T) :=
	\mathrm{Tr}(ST)$, with $S \in B_1(H)$ as per Definition
	\ref{def:trace-class}, is continuous, generalizing to the Banach space
	$B_1(H)$ the same weak$^*$ construction that Definition 1.6.28 performs
	on the topological dual of a $C^*$-algebra.
\end{definition}

\begin{remark}
	The ultraweak topology of Definition~\ref{def:ultraweak} is finer than
	the weak operator topology of Definition~\ref{def:wot-sot}, since the
	latter only tests continuity against rank-one operators $S =
	(\cdot,\phi)\psi$, a proper subset of $B_1(H)$. The two topologies
	coincide on norm-bounded subsets of $B(H)$, so that a von Neumann
	algebra $M$, as per Definition~\ref{def:vN-algebra}, closed in one is
	automatically closed in the other, consistently with
	Theorem~\ref{thm:bicommutant} already identifying both with the same
	purely algebraic condition. This equivalence is proved in \cite[TODO:
	collocazione precisa]{TODO-traceclass}.
\end{remark}

Restricting the ultraweakly continuous functionals of
Definition~\ref{def:ultraweak} from the whole of $B(H)$ to a von
Neumann subalgebra $M$ produces a distinguished Banach space attached
to $M$ itself, one that supplies the correct, structural, notion of a
normal state below.

\begin{definition}[Predual]\label{def:predual}
	Let $M$ be a von Neumann algebra on $H$, as per
	Definition~\ref{def:vN-algebra}. The predual of $M$ is the space
	\[
	M_* := \big\{ \omega|_M \mid \omega \in B_1(H) \big\}
	\]
	of restrictions to $M$ of the trace functionals $S \mapsto
	\mathrm{Tr}(S \, \cdot \,)$, as per Definition~\ref{def:trace-class},
	equipped with the quotient norm inherited from $B_1(H)$.
\end{definition}

\begin{remark}
	The predual $M_*$ of Definition~\ref{def:predual} is a Banach space, as
	per Definition 1.6.9, and the pairing $(\phi, A) \mapsto \phi(A)$, for
	$\phi \in M_*$ and $A \in M$, identifies $M$ isometrically with the
	topological dual $(M_*)'$ of $M_*$, as per Definition 2.2.13, and this
	identification is unique. Both facts are proved in \cite[TODO:
	collocazione precisa]{TODO-predual}.
\end{remark}

With the predual in hand, normality of a state can be expressed
intrinsically, as membership in $M_*$, rather than through the monotone
limits that only characterize it after the fact.

\begin{definition}[Normal state]\label{def:normal-state}
	Let $M$ be a von Neumann algebra on $H$, as per
	Definition~\ref{def:vN-algebra}. A state $\omega$ on $M$, as per
	Definition 2.1.3, is normal if $\omega \in M_*$, the predual of $M$ as
	per Definition~\ref{def:predual}.
\end{definition}

\begin{remark}
	Definition~\ref{def:normal-state} is equivalent to the more classical
	requirement that
	\[
	\omega\Big(\sup_\alpha A_\alpha\Big) = \sup_\alpha \omega(A_\alpha),
	\]
	for every bounded increasing net $(A_\alpha)$ of positive elements of
	$M$, with the supremum on the left taken in the strong operator
	topology of Definition~\ref{def:wot-sot}, and equally to continuity of
	$\omega$ for the ultraweak topology of Definition~\ref{def:ultraweak}
	restricted to $M$. This equivalence is proved in \cite[TODO:
	collocazione precisa]{TODO-normalStates}. When $M = B(H)$, it recovers
	the density-matrix description of Definition 2.3.15: normality of
	$\omega$ is equivalent to the existence of $\rho \in B_1(H)$ positive
	with $\mathrm{Tr}(\rho) = 1$ such that $\omega(A) = \mathrm{Tr}(\rho A)$
	for every $A \in B(H)$, matching Proposition 2.3.14.
\end{remark}

Every von Neumann algebra of Definition~\ref{def:vN-algebra} carries a
predual by construction, since Definition~\ref{def:predual} builds
$M_*$ directly from the ambient Hilbert space $H$. The converse
question, whether an abstract $C^*$-algebra with no reference to any
Hilbert space can already be recognized as a von Neumann algebra from
the existence of some predual alone, is the content of the abstract
characterization due to Sakai, and it is the form in which the
correspondence with Chapters 1 and 2 is used from Section 4.2 onward.

\begin{definition}[$W^*$-algebra]\label{def:wstar-algebra}
	A $W^*$-algebra is a $C^*$-algebra $A$, as per Definition 1.6.11, for
	which there exists a Banach space $A_*$, as per Definition 1.6.9, such
	that $A$ is isometrically isomorphic, as a Banach space, to the
	topological dual $(A_*)'$ of $A_*$, as per Definition 2.2.13.
\end{definition}

\begin{theorem}[Sakai]\label{thm:sakai}
	Let $A$ be a $W^*$-algebra, as per Definition~\ref{def:wstar-algebra}.
	There exist a Hilbert space $H$ and a $*$-isomorphism, as per
	Definition 2.3.1, of $A$ onto a von Neumann algebra $M$ on $H$, as per
	Definition~\ref{def:vN-algebra}. Moreover, the Banach space $A_*$ of
	Definition~\ref{def:wstar-algebra} is then unique, and this isomorphism
	identifies it with the predual $M_*$ of Definition~\ref{def:predual}.
\end{theorem}
\noindent This is proved in \cite[TODO: collocazione precisa]{TODO-Sakai}.

Theorem~\ref{thm:sakai} addresses the abstract side of the
correspondence, but it leaves open a question specific to the
construction already carried out in Chapter 2: nothing in the
Gelfand-Naimark-Segal construction of Theorem 2.7 was shown to produce
a weakly closed algebra of operators, so $\pi_\omega(A)$ was, until now,
only a $C^*$-subalgebra of $B(H_\omega)$, not yet known to be a von
Neumann algebra in the sense of Definition~\ref{def:vN-algebra}.
Normality of the generating state closes exactly this gap.

\begin{proposition}[GNS representations of normal states are weakly
	closed]\label{prop:gns-normal}
	Let $M$ be a von Neumann algebra on $H$, as per
	Definition~\ref{def:vN-algebra}, and let $\omega$ be a normal state on
	$M$, as per Definition~\ref{def:normal-state}. Let $(H_\omega,
	\pi_\omega, \Psi_\omega)$ be the Gelfand-Naimark-Segal representation
	of $\omega$, as per Theorem 2.7. Then $\pi_\omega(M)$ is a von Neumann
	algebra on $H_\omega$, as per Definition~\ref{def:vN-algebra}.
\end{proposition}
\noindent This is proved in \cite[TODO: collocazione precisa]{TODO-GNSnormal}.

\begin{remark}
	Proposition~\ref{prop:gns-normal} is what allows Sections 4.2 onward to
	treat $\pi_\omega(M)$ itself as a concrete von Neumann algebra, with
	$\Psi_\omega$ cyclic for it by part i) of Theorem 2.7, rather than
	merely as a $C^*$-subalgebra of $B(H_\omega)$ to which
	Theorem~\ref{thm:bicommutant} would not yet apply.
\end{remark}

The vector $\Psi_\omega$ obtained this way is cyclic by construction,
and its associated vector state is normal by the following general
fact, which holds for an arbitrary vector in a von Neumann algebra
without any reference to the Gelfand-Naimark-Segal construction itself.

\begin{remark}[Vector states are normal]\label{rem:vector-normal}
	Let $M$ be a von Neumann algebra on $H$, as per
	Definition~\ref{def:vN-algebra}, and let $\Psi \in H$ be a unit vector.
	Setting $\rho := P_\Psi$, the rank-one projector of Definition 2.2.18,
	the vector state $\omega(A) := (\Psi, A\Psi)$ coincides with the
	restriction to $M$ of the trace functional $A \mapsto \mathrm{Tr}(\rho
	A)$, as per Definition~\ref{def:trace-class}, since $\mathrm{Tr}(\rho
	A) = (\Psi, A\Psi)$ by the same computation as in Proposition 2.3.14.
	Hence every vector state on $M$ lies in the predual $M_*$ of
	Definition~\ref{def:predual}, and is in particular normal.
\end{remark}

Separation for $\pi_\omega(M)$ is a strictly stronger requirement than
cyclicity, and it is exactly the combination of both properties in a
single vector, on a von Neumann algebra rather than a bare
$C^*$-algebra, that Tomita's construction in Section 4.3 takes as its
point of departure.

\begin{lemma}[Cyclic and separating vectors]\label{lem:cyclic-separating}
	Let $M$ be a von Neumann algebra on $H$, as per
	Definition~\ref{def:vN-algebra}, and let $\Omega \in H$ be a unit
	vector, cyclic for $M$ as per Definition 2.2.34. Let $\omega$ be the
	associated vector state, $\omega(A) := (\Omega, A\Omega)$ for
	$A \in M$. Then $\Omega$ is separating for $M$, meaning that
	$A\Omega = 0$ with $A \in M$ forces $A = 0$, if and only if $\omega$ is
	faithful, as per Definition 2.3.17.
\end{lemma}

\begin{proof}
	Suppose $\omega$ is faithful and let $A \in M$ satisfy $A\Omega = 0$.
	Then $\omega(A^*A) = (\Omega, A^*A\Omega) = \lVert A\Omega \rVert^2 =
	0$, and faithfulness forces $A^*A = 0$. The $C^*$-identity of
	Definition 1.6.11 then yields $\lVert A \rVert^2 = \lVert A^*A \rVert =
	0$, that is, $A = 0$. Conversely, suppose $\omega$ is not faithful.
	There exists a nonzero $B \in M$ with $\omega(B^*B) = 0$, and the same
	computation gives $\lVert B\Omega \rVert^2 = \omega(B^*B) = 0$, so
	$B\Omega = 0$ with $B \neq 0$. Hence $\Omega$ fails to be separating
	for $M$.
\end{proof}

Remark 2.3.18 already showed that the Gelfand-Naimark-Segal
representation of a faithful state is itself faithful, by tracking the
kernel of the representation directly.
Lemma~\ref{lem:cyclic-separating} isolates the stronger statement
needed here, that the algebra acts on its own cyclic vector without
collapsing any nonzero element to zero, since it is this separating
property, rather than injectivity of the representation alone, that
Tomita's construction in Section 4.3 takes as its starting hypothesis.

\section{Factors, Projections, and the Trace Problem}\label{sec:factors}

Section 4.1 fixed the ambient algebraic setting for the rest of this
chapter, a von Neumann algebra $M$ together with a cyclic vector, and
the notion of a normal state guaranteeing that expectation values on
$M$ interact correctly with the operator topologies of
Definition~\ref{def:wot-sot}. Before constructing the modular flow
itself, it is worth understanding why a construction already available
in this thesis, representing a state through a density matrix as in
Definition 2.3.15 and Proposition 2.3.14, cannot serve as the general
framework this chapter needs. That construction rests on the trace of
Definition~\ref{def:trace} being available on the whole of $B(H)$, and
on the possibility of writing every normal state as $\omega(A) =
\mathrm{Tr}(\rho A)$ for some positive $\rho$ with $\mathrm{Tr}(\rho) =
1$. This section shows that such a trace need not exist on an
arbitrary von Neumann algebra at all, and classifies exactly when it
does, through the internal structure of the lattice of projections of
$M$. The classification singles out a family of algebras, of type III,
on which no nonzero trace of any kind is available, so that density
matrices in the sense of Chapter 2 simply do not exist there. It is
precisely this obstruction, rather than a technical inconvenience, that
motivates replacing the trace with the modular automorphism group
constructed in the remainder of the chapter, a substitute that requires
only a cyclic and separating vector, as isolated in
Lemma~\ref{lem:cyclic-separating}, and that reduces to the usual Gibbs
description of Section 4.6 whenever a trace happens to be available.

The comparison of projections that drives this classification is
carried out through operators that behave isometrically on part of $H$
and vanish on the rest, generalizing the unitary operators already
introduced in Definition 2.2.22 to the case where the initial and final
spaces need not coincide.

\begin{definition}[Partial isometry]\label{def:partial-isometry}
	Let $H$ be a Hilbert space, as per Definition 2.2.6. An operator
	$U \in B(H)$ is a partial isometry if $U^*U$ is a projector, as per
	Definition 2.2.17.
\end{definition}

\begin{remark}[Initial and final projections]\label{rem:partial-isometry}
	Let $U \in B(H)$ be a partial isometry, as per
	Definition~\ref{def:partial-isometry}, and write $P := U^*U$, the
	initial projection of $U$. For every $\psi \in H$,
	\[
	\lVert U\psi \rVert^2 = (\psi, U^*U\psi) = (\psi, P\psi) = \lVert
	P\psi \rVert^2,
	\]
	using $P^* = P = P^2$ as per Definition 2.2.17. Hence $\ker U = \ker
	P$, so that $U$ vanishes on $\ker P = (\operatorname{ran} P)^\perp$
	and restricts to an isometry on $\operatorname{ran}(P)$, a closed
	subspace by Theorem 2.4. Since $U$ vanishes on $\ker P$, it satisfies
	$UP = U$ as operators on $H$, so that $Q := UU^*$ satisfies
	\[
	Q^2 = UU^*UU^* = U(U^*U)U^* = UPU^* = UU^* = Q,
	\]
	while $Q^* = Q$ holds trivially. Hence $Q$ is itself a projector, as
	per Definition 2.2.17, called the final projection of $U$, with range
	equal to $\operatorname{ran}(U)$.
\end{remark}

Partial isometries supply exactly the notion of comparison needed to
declare two projections of a von Neumann algebra equal in size, even
when they are not literally the same operator, a comparison that
underlies the whole classification to follow.

\begin{definition}[Murray-von Neumann equivalence]\label{def:mv-equivalence}
	Let $M$ be a von Neumann algebra on $H$, as per
	Definition~\ref{def:vN-algebra}. Two projectors $P, Q \in M$, as per
	Definition 2.2.17, are Murray-von Neumann equivalent, written $P \sim
	Q$, if there exists a partial isometry $U \in M$, as per
	Definition~\ref{def:partial-isometry}, with $U^*U = P$ and $UU^* = Q$.
\end{definition}

\begin{remark}
	Definition~\ref{def:mv-equivalence} defines an equivalence relation on
	the projectors of $M$. Reflexivity holds since $P$ itself is a partial
	isometry with $P^*P = P^2 = P$, so $U = P$ witnesses $P \sim P$.
	Symmetry holds since $U$ witnessing $P \sim Q$ forces $U^*$, with
	$(U^*)^*U^* = UU^* = Q$ and $U^*(U^*)^* = U^*U = P$, to witness $Q \sim
	P$. For transitivity, let $U_1$ witness $P \sim Q$ and $U_2$ witness $Q
	\sim R$. Since $Q$ is the final projection of $U_1$ and the initial
	projection of $U_2$, Remark~\ref{rem:partial-isometry} yields $QU_1 =
	U_1$ and $U_2Q = U_2$, so that $U := U_2U_1 \in M$ satisfies
	\[
	U^*U = U_1^*U_2^*U_2U_1 = U_1^*QU_1 = (QU_1)^*U_1 = U_1^*U_1 = P,
	\qquad UU^* = U_2U_1U_1^*U_2^* = U_2QU_2^* = U_2U_2^* = R,
	\]
	so $U$ witnesses $P \sim R$.
\end{remark}

Murray-von Neumann equivalence compares projections without reference
to any external notion of size, but the classification of factors also
needs to compare a projection with a strictly smaller one contained
inside it, a containment already implicit in the range inclusions used
above.

\begin{remark}[Order on projectors]\label{rem:projector-order}
	For projectors $P, Q \in M$, as per Definition 2.2.17, write $P \leq
	Q$ when $PQ = P$. This holds if and only if $\operatorname{ran}(P)
	\subseteq \operatorname{ran}(Q)$, and it turns the set of projectors of
	$M$ into a partially ordered set.
\end{remark}

With both a notion of equivalence and a notion of containment
available, a projection can be tested against its own equivalents
lying strictly inside it, and the outcome of this test is what
separates the finite from the infinite case in the trichotomy below.

\begin{definition}[Finite and infinite projections]\label{def:finite-projection}
	Let $M$ be a von Neumann algebra on $H$, as per
	Definition~\ref{def:vN-algebra}, and let $P \in M$ be a projector, as
	per Definition 2.2.17. The projector $P$ is finite if every projector
	$Q \in M$ satisfying $Q \leq P$, as per Remark~\ref{rem:projector-order},
	and $Q \sim P$, as per Definition~\ref{def:mv-equivalence}, satisfies
	$Q = P$. The projector $P$ is infinite if it is not finite.
\end{definition}

Having compared projections against each other, the classification
turns to the algebra as a whole, through the elements that commute
with every one of its own members.

\begin{definition}[Center]\label{def:center}
	Let $M$ be a von Neumann algebra on $H$, as per
	Definition~\ref{def:vN-algebra}. The center of $M$ is $Z(M) := M \cap
	M'$, with $M'$ the commutant of $M$, as per Definition~\ref{def:commutant}.
\end{definition}

\begin{remark}
	The center $Z(M)$ of Definition~\ref{def:center} is itself abelian.
	Indeed, for $A, B \in Z(M)$, $A \in M$ and $B \in Z(M) \subseteq M'$
	force $AB = BA$, since every element of $M'$ commutes with every
	element of $M$ by Definition~\ref{def:commutant}.
\end{remark}

\begin{definition}[Factor]\label{def:factor}
	A von Neumann algebra $M$ on $H$, as per Definition~\ref{def:vN-algebra},
	is a factor if $Z(M) = \mathbb{C}1$, with $Z(M)$ the center of $M$, as
	per Definition~\ref{def:center}.
\end{definition}

The factors are the von Neumann algebras with no nontrivial invariant
splitting into independent pieces, and every von Neumann algebra
decomposes into a direct integral of factors. The trichotomy below
concerns exactly these building blocks. Its statement, however,
requires generalizing the finite-dimensional trace already used
informally in Definition 2.3.15 to an object that need not be bounded
and need not even be finite, so that it can register the difference
between a factor on which such a functional exists and one on which it
does not.

\begin{definition}[Trace]\label{def:trace-functional}
	Let $M$ be a von Neumann algebra on $H$, as per
	Definition~\ref{def:vN-algebra}, and let $M_+$ denote its positive
	elements, as per Definition 2.2.22. A trace on $M$ is a map $\tau :
	M_+ \to [0, \infty]$ satisfying
	\[
	\tau(A + B) = \tau(A) + \tau(B), \qquad \tau(\lambda A) = \lambda
	\tau(A), \qquad \tau(UAU^*) = \tau(A),
	\]
	for all $A, B \in M_+$, all $\lambda \geq 0$, with the convention $0
	\cdot \infty := 0$, and all unitary $U \in M$, as per Definition 2.2.22.
\end{definition}

\begin{definition}[Faithful trace]\label{def:faithful-trace}
	Let $\tau$ be a trace on $M$, as per Definition~\ref{def:trace-functional}.
	The trace $\tau$ is faithful if $\tau(A) = 0$ with $A \in M_+$ forces
	$A = 0$.
\end{definition}

\begin{definition}[Normal trace]\label{def:normal-trace}
	Let $\tau$ be a trace on $M$, as per Definition~\ref{def:trace-functional}.
	The trace $\tau$ is normal if
	\[
	\tau\Big(\sup_\alpha A_\alpha\Big) = \sup_\alpha \tau(A_\alpha),
	\]
	for every bounded increasing net $(A_\alpha)$ of elements of $M_+$,
	with the supremum on the left taken in the strong operator topology of
	Definition~\ref{def:wot-sot}.
\end{definition}

\begin{definition}[Semifinite trace]\label{def:semifinite-trace}
	Let $\tau$ be a trace on $M$, as per Definition~\ref{def:trace-functional}.
	The trace $\tau$ is semifinite if, for every nonzero $A \in M_+$, there
	exists a nonzero $B \in M_+$ with $B \leq A$, as per
	Remark~\ref{rem:projector-order} extended to the order induced on all
	of $M_+$, and $\tau(B) < \infty$.
\end{definition}

\begin{remark}
	Definition~\ref{def:normal-trace} is the same continuity condition as
	Definition~\ref{def:normal-state}, extended from normalized states to
	the possibly unbounded functional $\tau$. When $M = B(H)$ and $\tau =
	\mathrm{Tr}$, as per Definition~\ref{def:trace}, restricted to the
	whole of $B(H)_+$ rather than to $B_1(H)_+$, $\tau$ is faithful, normal,
	and semifinite, and every density matrix $\rho$ of Definition 2.3.15
	satisfies $\mathrm{Tr}(\rho) < \infty$, consistently with
	Definition~\ref{def:semifinite-trace} applied to $A = 1$.
\end{remark}

With finiteness of projections, factoriality, and the trace all in
place, the classification of Murray and von Neumann can now be stated.
It shows that the existence of a trace on a factor is entirely
determined by the combinatorics of its projections, so that the
absence of minimal projections already forces a choice between two
mutually exclusive possibilities for the whole algebra.

\begin{theorem}[Murray-von Neumann classification of factors]\label{thm:murray-von-neumann}
	Let $M$ be a factor on $H$, as per Definition~\ref{def:factor}. Exactly
	one of the following holds.
	\begin{enumerate}
		\item Type I: $M$ has a minimal nonzero projector, that is, a
		projector $P \in M$, as per Definition 2.2.17, such that no projector
		$Q \in M$ satisfies $0 \neq Q < P$ as per
		Remark~\ref{rem:projector-order}. In this case $M$ is $*$-isomorphic,
		as per Definition 2.3.1, to $B(K)$ for some Hilbert space $K$.
		\item Type II: $M$ has no minimal nonzero projector, but $M$ admits a
		faithful normal semifinite trace, as per Definitions
		\ref{def:trace-functional}, \ref{def:faithful-trace},
		\ref{def:normal-trace}, and \ref{def:semifinite-trace}.
		\item Type III: $M$ has no minimal nonzero projector, and $M$ admits
		no nonzero faithful normal semifinite trace.
	\end{enumerate}
\end{theorem}
\noindent This is proved in \cite[TODO: collocazione precisa]{TODO-MurrayVonNeumann}.

\begin{remark}
	Type III factors are not a single, further indistinguishable class.
	This third case of Theorem~\ref{thm:murray-von-neumann} splits into a
	family of mutually exclusive subtypes, denoted $\mathrm{III}_0$,
	$\mathrm{III}_\lambda$ for $\lambda \in (0,1)$, and $\mathrm{III}_1$,
	detected by an invariant built from the modular operator only once the
	tools of Section 4.4 are available, and Section 4.9 returns to this
	finer classification in full. \cite[TODO: collocazione
	precisa]{TODO-fineTypeIII}
\end{remark}

The type III case of Theorem~\ref{thm:murray-von-neumann} is where the
framework of Chapter 2 breaks down completely, and it is worth spelling
out exactly what is lost.

\begin{remark}[Collapse of the density matrix picture]\label{rem:type-III-density-matrix}
	Let $M$ be a factor of type III on $H$, as per
	Theorem~\ref{thm:murray-von-neumann}. Since $M$ admits no nonzero
	faithful normal semifinite trace, no normal state $\omega$ on $M$, as
	per Definition~\ref{def:normal-state}, can be represented as
	$\omega(A) = \tau(\rho A)$ for a fixed positive $\rho \in M$ with
	respect to any trace $\tau$ on $M$, in contrast with the description of
	Definition 2.3.15 and Proposition 2.3.14 available for $B(H)$ itself.
	Every normal state on $M$ is nevertheless still a vector state for some
	faithful representation of $M$, by Remark~\ref{rem:vector-normal}
	together with the Gelfand-Naimark-Segal construction of Theorem 2.7,
	so a Hilbert space description of the state survives even though a
	density matrix does not. Algebras of local observables in relativistic
	quantum field theory are type III on physical grounds, so that this
	collapse is the generic situation for such algebras rather than a
	pathological exception. \cite[TODO: collocazione precisa]{TODO-QFTtypeIII}
\end{remark}

Theorem~\ref{thm:murray-von-neumann} therefore leaves this thesis with
a genuine gap. Chapter 2 built the entire probabilistic content of
quantum mechanics, states as density matrices and time evolution as
conjugation by $e^{-itH/\hbar}$, on the trace of $B(H)$, a type I
factor by Theorem~\ref{thm:murray-von-neumann} applied with $M = B(H)$
itself, since every rank-one projector is minimal there. On a type III
factor, neither notion is available in that form. The remainder of
this chapter constructs, for an arbitrary von Neumann algebra $M$
together with a vector cyclic and separating for it as in
Lemma~\ref{lem:cyclic-separating}, a canonical one-parameter group of
automorphisms of $M$ that requires no trace whatsoever, built instead
directly from the vector state itself through Tomita's theorem of
Section 4.4. Section 4.6 will verify explicitly that, when $M$ does
possess a trace, this group reduces to ordinary Gibbs time evolution,
so that the construction to follow genuinely extends, rather than
replaces, the familiar type I picture of Chapter 2.
\end{document}